%%%%%%%%%%%%%%%%%%%%%%%%%%%%%%%%%%%%%%%%%%%%%%%%%%%%%%%%%%%%%%%%%%%%%%%%%%%%%%%%%%%%%%%%%%%%%%%%%
% Chapter 6: Discussion   (the \chapter command is in main.tex)
%%%%%%%%%%%%%%%%%%%%%%%%%%%%%%%%%%%%%%%%%%%%%%%%%%%%%%%%%%%%%%%%%%%%%%%%%%%%%%%%%%%%%%%%%%%%%%%%%

Chapter~\ref{ch:results} showed three main things. An EKF that learns the cell resistance tracks the SoC far better than charge counting. The five-state filter is the only one whose uncertainty stays honest across healthy, biased and aged batteries. And a decision rule that uses that uncertainty keeps its safety promise while a rule that ignores it does not. This chapter explains why these results occur, what limits them, and what would improve them.

\section{Why the EKF works for this problem}\label{sec:d_why_ekf}

The EKF works here because the voltage gives it an independent check on the SoC at every second, and the Kalman gain decides how much to trust that check (Table~\ref{tab:bg_gain}). Charge counting has no such check. In E1 its error grew steadily from its start error through the sensor bias, to $-5.50\%$ at landing, while stage~B stayed within 1.53\% (RMSE 0.97\% against 4.36\%). In E2 charge counting kept almost all of its 15-point start error (RMSE 13.85\%), while stage~B removed it within 42\,s (RMSE 1.32\%). The covariance did what it should: it started wide when the start was uncertain, let the gain be large, and shrank as the voltage supplied information.

The tangent-line approximation was adequate. The OCV curve is smooth, and even in E2, with a 20\% initial uncertainty and a 15-point start error, the filter converged without any sign of divergence. The design-stage ablation (Table~\ref{tab:m_ablation}) gives a sharper test: with a perfect model and no bias, the two-state filter reached an SoC RMSE of 0.12\%. Linearisation errors were therefore small; the errors that remained in every later experiment came from model mismatch, not from the EKF's approximation. Finally, the noise in the simulation is Gaussian and white by construction, which matches the EKF's assumptions; real sensors may be less kind (Section~\ref{sec:d_limitations}).

\section{Gain, OCV slope and uncertainty}\label{sec:d_gain_slope}

The OCV slope sets how much the voltage can say about the SoC, because it is the first entry of $H$ \eqref{eq:m_H}. Where the curve is steep, a small SoC error gives a large voltage error, the gain is large and the SoC uncertainty stays small; where it is flat, the opposite holds. E7 showed this directly for stage~C: its $\sigma_{\text{SoC}}$ rose from about 1.5\% to about 2.2\% as the slope fell from about 8 to about 5\,mV per 1\% SoC, and it dipped briefly exactly where the slope has a local bump near 65--70\% SoC. Because the uncertainty follows the slope rather than simply the time since take-off, the effect is caused by the information in the voltage.

This matters for the decision rule. The middle of the SoC range, where the curve is flat, is where the turn-back decisions of the long missions are made, and there stage~C reports a larger $\sigma_{\text{SoC}}$. Policy~D's safety buffer, $1.645\,\sigma_{\text{total}}$, grows accordingly. The drone becomes more cautious exactly where the voltage is a weaker clue.

Stage~B behaved differently: its uncertainty fell steadily from about 0.95\% to about 0.70\% through the flight. Every change in current gives it more information about $R_0$, and better knowledge of $R_0$ reduces the SoC uncertainty, which masks the flat-region effect. The stated uncertainty therefore depends on the structure of the whole filter, not on the OCV curve alone. The same slope effect also explains a detail of the experiments: near full charge the curve rises by about 10\,mV per 1\% SoC, so a single noisy resting reading (10\,mV noise) gives a starting SoC uncertain by about 1\%. The start about 3\% low in the first E1 run was an unlucky case of this, and averaging 10 readings reduced it.

\section{Observability of resistance, bias and capacity}\label{sec:d_observability}

Whether a state can be learned depends on whether the measurements contain information that separates it from the others. The resistance $R_0$ is a good example of a state that needs \emph{excitation}. During the constant-current take-off, a voltage that is lower than expected can be explained equally well by a lower SoC or by a higher $R_0$. The SoC entry of $H$, $d\,\mathrm{OCV}/d\,\mathrm{SoC}$, does not depend on the current, while the $R_0$ entry, $-(I_m - \hat{b})$, does. When the current stepped from 30\,A to 15\,A at about 0.5\,min in E1, the two became separable: within seconds the stage-B SoC error moved from about $-3\%$ to about $+1\%$, its $\pm 3\sigma$ band shrank from about $\pm 8.5\%$ to about $\pm 2.5\%$, and $R_0$ moved from about 8\,m$\Omega$ to about 8.9\,m$\Omega$, close to the true 8.8\,m$\Omega$. The same mechanism delayed stage~B's recovery in E2, where it stayed 5--7\% too high until the step. A delivery drone gets this excitation naturally from take-off, hover and cruise; a flight at perfectly constant current would not.

The bias and the capacity are harder. Equation~\eqref{eq:m_confound} shows that, at a steady current, a bias error and a resistance error change the voltage in the same way, and the capacity does not appear in the voltage at all. Both affect the SoC only through the discharge rate $(I_m - b)/C_n$ \eqref{eq:m_rate}, so within one flight only that rate is observable. The experiments show both sides of this. Stage~C did not pin down the bias in E3 (it ended at $+0.33$\,A against $+0.80$\,A, with a band of about $\pm 1.8$\,A), and in E4 it ended with a wrong-signed bias ($-0.42$\,A against $+0.50$\,A) and a capacity of 4.45\,Ah against 4.25\,Ah. Yet in E4 its rate was within about 1.5\% of the truth and its final SoC error was only $+0.17\%$.

The voltage feedback also explains why an unlearned bias did little harm. In the design-stage ablation a 0.5\,A bias alone gave stage~A an RMSE of only 0.43\% (Table~\ref{tab:m_ablation}), and in E3 stage~B, which has no bias state, reached 1.28\% against 3.13\% for charge counting. The drift that a bias causes is slow, and the voltage corrects it continuously. An unlearned capacity error is different, because it can hide behind $R_0$, as E4 showed for stage~B (Section~\ref{sec:d_choice}).

\section{Effect of model mismatch}\label{sec:d_mismatch}

A Kalman filter assumes that its errors are white noise. Model mismatch breaks this assumption, because it produces errors that have the same sign for minutes at a time. The design-stage investigation showed the consequence clearly: making the filter trust the voltage less by inflating $R$ did not reduce the error of the two-state filter (RMSE 2.49\%), because the systematic resistance error still pointed it to the wrong SoC, only more slowly. Learning the dominant error, $R_0$, did reduce it. The 10\,mV model-error allowance then serves a different purpose: once the large error has been removed by learning $R_0$, it widens the bands enough to cover the smaller errors that remain. The cost is a mean NIS of about 0.5, a slightly cautious filter; the gain is that stage~C's band contained the truth 100\% of the time in E1--E4.

Model mismatch is also hard to detect from the filter's own signals. Stage~A had a mean NIS of about 0.5 in every healthy scenario, which looks healthy, while its mean NEES was between 83 and 93 in E1--E3 and reached 404 in E6 (Table~\ref{tab:r_e8_extra}). Stage~A explained the voltage well, but with the wrong SoC. Innovation-based checks test whether the measurements are explained, not whether the hidden state is right. On a real drone, where the true SoC is unknown, this kind of failure could not be seen from the innovations alone, which is a strong argument for testing a filter in simulation against a deliberately different truth model, as this project does.

The remaining mismatch in the default true battery is the slow RC pair. It caused the gap between OCV and terminal voltage to widen slowly in the simulator verification (Section~\ref{sec:r_sim}), and it is the most likely cause of stage~B's slow drift to $-1.53\%$ in E1, which stayed inside its band. E6 showed the limit when this mismatch is made large ($R_2 = 4$\,m$\Omega$). Every filter then drifted low; stage~A ended at $-7.91\%$, stage~C left its band around 12\,minutes, and no filter's band held the truth more than 83\% of the time over 20 flights. Charge counting, which ignores the voltage and is not affected by the slow RC pair, was as accurate as the best filter. When the model is too wrong, the voltage feedback misleads the EKF. The honest remedies are a better model (a second RC state in the filter), a larger model-error allowance, or adaptive noise estimation that raises $R$ when the innovations show a trend.

\section{Tuning}\label{sec:d_tuning}

E5 showed that the noise settings trade accuracy against honesty, and that only one kind of mistuning is visible from the filter's own signals. With $R{\times}0.01$ the filter believed every wiggle of the voltage, swung to about $-13\%$ during take-off before $R_0$ was learned, and was badly overconfident (30\% coverage). Its mean NIS of 32.34 made the mistake obvious. With $R{\times}100$ the band became far too wide (NIS 0.01), which would make the RTH rule turn back needlessly early. With $Q{\times}100$ the estimate became jumpy (RMSE 2.35\%). The most instructive case was $Q{\times}0.01$: it was the most accurate in its single run (RMSE 0.67\% against 1.06\%), but its band held the truth only 94\% of the time, because a filter that trusts its model more has less room for the model errors it does not know about. The chosen tuning was kept, accepting slightly lower accuracy for a band that stayed honest, because the decision rule turns $\sigma_{\text{SoC}}$ directly into a safety margin. Together with the stage-A result of Section~\ref{sec:d_mismatch}, E5 shows what NIS can and cannot do: it detects a wrongly set $R$, but not a biased model.

\section{Choosing the filter for decisions}\label{sec:d_choice}

The usual way to choose between filters is accuracy: the one with the lowest error wins. By that measure stage~B is the best filter in this project. On a healthy battery its mean SoC RMSE over 20 flights was 1.03\%, against 1.41\% for stage~C (Table~\ref{tab:r_e8}). The aged-battery experiment E4 shows why accuracy alone is the wrong test for a filter that feeds a safety decision. When the capacity is lower than assumed, the voltage falls faster than stage~B expects. During steady cruise a lower voltage can be explained either by a lower SoC or by a higher resistance, and stage~B explains it with its $R_0$ state. Its SoC estimate therefore stays too high: it ended on average 8.34\% above the truth, its band contained the truth only 27\% of the time, and its mean NEES was 45.6. It was not only wrong but confidently wrong, and on the optimistic side, which is the side that makes a drone turn back too late.

Stage~C has more unknowns to share the same information, so it is a little less accurate on a healthy battery. In exchange it was the only filter whose band contained the truth 100\% of the time in all four scenarios E1--E4. It achieves this although its individual parameter estimates are often wrong. In E4 it ended with a capacity of 4.45\,Ah (true 4.25\,Ah) and a bias of $-0.42$\,A (true $+0.50$\,A), yet its final SoC error was $+0.17\%$. The reason is the observability structure of the model (Section~\ref{sec:m_ekf_choice}): the bias and the capacity affect the SoC only through the discharge rate $(I_m - b)/C_n$, and the voltage checks only that rate. Stage~C got the rate right to within about 1.5\% and was honest about not knowing the two parts separately; in E3 its bias estimate kept a band of about $\pm 1.8$\,A.

The Monte Carlo study confirmed the choice in the setting that matters. Over 500 long missions with random battery ages, biases and parameters, stage~C had a mean RMSE of 1.31\% and a band coverage of 100\%, and the decision rule built on it had no deaths. The lesson generalises beyond batteries: \emph{the most accurate filter is not always the safest one to make decisions with}. A filter that feeds a decision should be judged by the honesty of its uncertainty across all the conditions it will meet, not only by its accuracy in the nominal case. The design also points to how stage~C could become more accurate without losing honesty: calibrating the current-sensor bias at rest before take-off, when the true current is known to be zero, and estimating the capacity between flights from the charge drawn and the SoC read from the resting voltage before and after. Both would make $b$ and $C_n$ individually observable.

\section{Value of uncertainty in the return-to-home decision}\label{sec:d_value_uncertainty}

Policies C and D are the clearest experiment in the project, because they share everything except the use of the uncertainty. They use the same stage-C filter, the same trip cost and the same reserve. The only difference is that C turns back when the estimated margin reaches zero, while D turns back when the probability of keeping the reserve falls below 95\%. Over 500 missions neither had a single death, which shows that the estimator itself was good. But C broke the reserve in 38\% of missions, and D in 2.8\%.

The reason is simple. C aims at the reserve. Its estimate and its trip cost are close to the truth, but never exact, so whenever they err slightly on the optimistic side the drone lands slightly below the reserve. If the errors were symmetric, this would happen in about half of the missions; 38\% is of that order. D aims above the reserve by $1.645\,\sigma_{\text{total}}$, a buffer that is large when the filter is unsure and small when it is confident. In the demonstration mission this buffer cost half a minute of flight and gave 7.8 points of SoC margin at landing instead of C's 1.9.

D's breach rates also show that the probabilities mean what they say. At every threshold in the sweep the share of missions that broke the reserve stayed within the promised $1 - p^{\ast}$, and at 80--95\% it was about half of the promise (Table~\ref{tab:r_sweep}). This is called \emph{calibration}. It is a test of the whole chain: the EKF's covariance, the trip cost, the trip allowance and the Gaussian assumption. The slight conservatism agrees with what E8 showed about stage~C (mean NIS about 0.5, mean NEES below 1), which comes from the deliberate 10\,mV model-error allowance. In a safety system, a promise that is kept with some room to spare is preferable to one that is broken. An honest uncertainty is what turns a filter into something a decision rule can rely on, and it is why the stage-C choice of Section~\ref{sec:d_choice} matters.

The two baselines show what happens without a reliable estimate. Charge counting (B) has no uncertainty and no way of noticing that it is wrong. With an aged battery or a sensor that reads low it is optimistic, and it died in 30.4\% of all missions and in 58.2\% of those with an aged battery. Its higher delivery count (1.56) is the result of recklessness, not of better use of the battery: on average it landed below the reserve. The voltage threshold (A) fails differently. The voltage at which a drone must turn back depends on how far it is from home and on the state of the battery, so no fixed value can be right for every mission. In the sweep, the thresholds that keep A safe waste a great deal of battery: at 3.55\,V it broke the reserve about as rarely as D at 95\% (3.0\% against 2.0\%), but left 17.9\% of the battery unused, about three times D's 6.1\%, and completed about half as many deliveries (0.56 against 1.18). Across the whole sweep, D's trade-off curve lies below A's: for every voltage threshold that keeps the breach rate below that of C, a D threshold is at least as safe and less wasteful.

Finally, the threshold sweep turns the choice of $p^{\ast}$ into an explicit trade-off. Moving from 90\% to 99\% raised the unused battery from 5.0\% to 8.2\% and reduced the deliveries from 1.22 to 1.15 per mission; moving on to 99.9\% cost a further 2.2 points of battery. An operator can read off the price of each extra ``nine'' and choose. The 95\% used here is a reasonable default: no deaths, a breach rate of about 2--3\%, and about 6\% of the battery left unused.

\section{Limitations}\label{sec:d_limitations}

The results should be read with the following limitations in mind.
\begin{itemize}
  \item \textbf{Simulation only.} All results come from simulation, where the true SoC is known and every error can be measured. The sensor noise is Gaussian and white by construction. Real flight logs have no true SoC, so the consistency checks that revealed stage~A's bias (NEES and band coverage) would not be available.
  \item \textbf{A representative OCV curve.} The OCV table has the typical shape of an NMC cell but was not measured on a specific cell. A real application would need the curve of its own cells, measured at rest; errors in it act like the 5\,mV offset tested here, but could be larger.
  \item \textbf{No temperature.} A cold battery has a higher resistance and less usable capacity. Learning $R_0$ would follow part of this, but the capacity change would face the same observability limits as ageing.
  \item \textbf{No hysteresis.} Lithium-ion cells show a small difference in resting voltage after charge and after discharge, which is ignored in both models.
  \item \textbf{Weakly observable bias and capacity.} Within one flight only their combined effect on the discharge rate is learned (E3, E4).
  \item \textbf{Strong model mismatch.} With a large unmodelled slow RC pair (E6) no filter was consistent, and the EKF was no better than charge counting.
  \item \textbf{Gaussian assumption in the decision rule.} $P_{\text{safe}}$ assumes Gaussian errors. Heavy-tailed events, such as a sudden gust or a motor fault, are not covered.
  \item \textbf{Simple trip model.} The cost of getting home assumes a straight path and a cruise current that stays at its learned value. Obstacles, no-fly zones or a change in wind on the way home are covered only by the 10\% trip allowance, and the uncertainty of the drain rate itself is not added to $\sigma_{\text{total}}$.
  \item \textbf{Landing-phase artefact.} During the final landing the trip cost still includes the full 40\,s landing, so $P_{\text{safe}}$ drops there (Section~\ref{sec:r_single_rth}). Decisions are not made in that phase, so the results are not affected, but the displayed probability is misleading.
  \item \textbf{One mission geometry.} All return-to-home missions used the same long route, with the same pattern of take-off, cruise and hover. Other routes would give different current changes, and so different excitation for learning $R_0$.
  \item \textbf{Simplifications in the filter.} The correlation between the current noise in the process and in the measurement is ignored, and the current-noise term of $R$ uses the nominal $R_0$.
\end{itemize}

\section{When a UKF or particle filter would be better}\label{sec:d_alternatives}

An unscented Kalman filter avoids the single tangent line by propagating a small set of sample points through the nonlinear functions, and a particle filter represents the uncertainty with many samples and can handle non-Gaussian errors~\cite{simon2006optimal}. Dibbern et al.\ found a UKF more accurate than an EKF for a multirotor UAV battery~\cite{purdue_jate_comparison}. In this project, however, the evidence points to model mismatch, not linearisation, as the main source of error: with a perfect model the EKF error was 0.12\%, and the errors in E1--E6 came from resistance, slow-RC and capacity differences. A UKF uses the same model and would not remove those errors. It would become more attractive with sharper nonlinearity, for example deep into the steep region near empty, with cell chemistries whose OCV curve has a very flat plateau, or with a much larger initial uncertainty. A particle filter would be justified mainly if the noise were strongly non-Gaussian or the uncertainty multimodal. A direct comparison on the same simulator would settle the question and is left as future work.

\section{Future work}\label{sec:d_future}

Several extensions follow directly from the results.
\begin{itemize}
  \item \textbf{Bias calibration at rest.} Before take-off the true current is known to be zero, so the current sensor's bias can be measured directly. This would remove one of the confounded unknowns of stage~C.
  \item \textbf{Capacity estimation between flights.} Comparing the charge drawn in a flight (ampere-hours) with the change in SoC read from the resting voltage before and after gives the capacity directly, without the single-flight observability problem.
  \item \textbf{A two-RC filter model} to remove the slow-RC mismatch that limited the filters in E1 and E6.
  \item \textbf{Adaptive noise estimation}, raising $R$ or $Q$ when the innovations show a trend, to keep the bands honest under a strong mismatch.
  \item \textbf{A UKF and particle-filter comparison} on the same simulator and missions.
  \item \textbf{Validation on real flight logs}, for example the CMU in-flight energy data of a delivery quadcopter~\cite{cmu_dji_dataset} and the eVTOL battery dataset~\cite{cmu_evtol_dataset,cmu_evtol_paper}, using consistency checks that do not need the true SoC.
  \item \textbf{A richer decision model}: temperature, varied routes, and the uncertainty of the drain rate added to $\sigma_{\text{total}}$.
\end{itemize}
